\documentclass[11pt]{article}
\usepackage[margin=1in]{geometry}
\usepackage{amsmath,amssymb,amsthm,mathtools,bm}
\usepackage[numbers,sort&compress]{natbib}
\usepackage{hyperref}
\usepackage{enumitem}
\usepackage{booktabs}

\hypersetup{
  colorlinks=true,
  linkcolor=blue,
  citecolor=blue,
  urlcolor=blue
}

\newcommand{\F}{\mathbb{F}_2}
\newcommand{\Prb}{\mathbb{P}}
\newcommand{\E}{\mathbb{E}}
\newcommand{\Var}{\mathrm{Var}}
\newcommand{\ind}[1]{\mathbf{1}\!\left[#1\right]}

\title{LiDMaS+ Lecture Notes:\\From Physical Noise to Threshold Inference}
\author{
Dennis Delali Kwesi Wayo\\
\href{https://orcid.org/}{ORCID: https://orcid.org/}
}
\date{\today}

\begin{document}
\maketitle
\tableofcontents

\section{Purpose and Scope}
LiDMaS+ is a simulator-oriented framework for threshold studies in planar surface codes, with both Pauli and hybrid CV-discrete noise paths \cite{wayo2026lidmas}. These notes explain how the implemented mathematical expressions are derived and why each expression appears in the pipeline.

The teaching goal is to make the equations ``narrative'':
\begin{enumerate}[label=\arabic*.]
  \item represent physical errors in algebraic form,
  \item infer likely causes from syndromes using a decoder objective,
  \item estimate logical failure rates and threshold indicators statistically.
\end{enumerate}

\paragraph{How to read these notes (self-study mode).}
For each section, ask three questions:
\begin{enumerate}[label=\arabic*.]
  \item What object is being modeled?
  \item What optimization or estimator is being applied?
  \item What quantity is exported to threshold plots?
\end{enumerate}

\section{Algebraic Backbone Over \texorpdfstring{$\F$}{F2}}
Pauli error presence/absence is binary, so the natural field is
\begin{equation}
\F = \{0,1\},\qquad a\oplus b=(a+b)\bmod 2.
\end{equation}
This choice ensures parity constraints and syndrome equations are exact modulo-2 statements.

For $H\in\F^{m\times n}$ and $e\in\F^n$:
\begin{equation}
(He)_i=\sum_j H_{ij}e_j \bmod 2.
\end{equation}
If $(He)_i=1$, check $i$ detects odd parity among incident error bits.

\paragraph{Variables used throughout.}
\begin{equation}
d:\text{code distance},\quad n:\text{data qubits},\quad m_X,m_Z:\text{check counts},
\end{equation}
\begin{equation}
p:\text{Pauli rate},\quad \sigma:\text{CV Gaussian scale},\quad N:\text{trials},\quad k:\text{failures}.
\end{equation}

\section{Surface-Code Geometry in the LiDMaS+ Implementation}
The surface-code setting follows the standard topological framework \cite{surfacecode1998}. For odd $d\ge 3$, LiDMaS+ uses:
\begin{equation}
n=2d(d-1),\qquad m_X=d^2,\qquad m_Z=(d-1)^2.
\end{equation}
Parity-check matrices are
\begin{equation}
H_X\in\F^{m_X\times n},\qquad H_Z\in\F^{m_Z\times n}.
\end{equation}
Given X- and Z-sector data errors $e_X,e_Z\in\F^n$:
\begin{equation}
s_Z = H_Z e_X \bmod 2,\qquad s_X = H_X e_Z \bmod 2.
\end{equation}

\paragraph{Interpretation.}
This cross-coupling is physically expected: Z-type stabilizers are sensitive to X-type faults and vice versa.

\section{Logical Predicate and Failure Variable}
Let logical supports be $L_X,L_Z\in\F^n$ and define
\begin{equation}
\langle a,b\rangle=\sum_i a_ib_i\bmod 2.
\end{equation}
After correction $c$, the residual is $e^{\mathrm{res}}=e\oplus c$.

In the current threshold harness, the implementation-aligned logical event is
\begin{equation}
\mathcal{L}=
\ind{
\langle e_X\oplus c,L_X\rangle=1
\ \lor\
\langle e_X\oplus c,L_Z\rangle=1
}.
\end{equation}
Hence $\mathcal{L}=1$ marks one logical failure trial.

\paragraph{Why this matters.}
All downstream statistics use $\mathcal{L}$ as the Bernoulli outcome. If you change this predicate, every reported threshold curve changes.

\section{Noise Modeling: Pauli and Hybrid CV Paths}
\subsection{Discrete Pauli mode}
Independent Bernoulli draws:
\begin{equation}
\Prb(e_{X,i}=1)=p,\qquad \Prb(e_{Z,i}=1)=p.
\end{equation}
For the active Z-syndrome threshold path, $e_X$ drives
\begin{equation}
s_Z=H_Ze_X\bmod 2.
\end{equation}

\subsection{Hybrid mode with GKP-style digitization}
Hybrid CV motivation follows GKP-style lattice quantization ideas \cite{gkp2001}:
\begin{equation}
\Delta q\sim\mathcal N(0,\sigma^2),\qquad
\Delta p\sim\mathcal N(0,\sigma^2).
\end{equation}
Using $\lambda=\sqrt{\pi}$:
\begin{equation}
n_q=\mathrm{round}\!\left(\frac{\Delta q}{\lambda}\right),\qquad
n_p=\mathrm{round}\!\left(\frac{\Delta p}{\lambda}\right).
\end{equation}
Binary flips are parity bits:
\begin{equation}
x_{\mathrm{flip}}=|n_q|\bmod 2,\qquad z_{\mathrm{flip}}=|n_p|\bmod 2.
\end{equation}
These flips define effective $(e_X,e_Z)$ passed to the same syndrome/decoder pipeline.

\paragraph{Self-check.}
The key conceptual move is a \emph{continuous-to-discrete map}. Once digitized, the decoding algebra is identical to Pauli mode.

\section{Decoder Objective I: MWPM}
Matching-based decoding is a standard reference strategy in surface-code benchmarking \cite{pymatching2021}.

\subsection{Defect graph and costs}
Defects are checks with syndrome value 1. For defects $i,j$:
\begin{equation}
d_{ij}=|x_i-x_j|+|y_i-y_j|.
\end{equation}
Uniform mode:
\begin{equation}
W_{ij}=d_{ij}.
\end{equation}
Weighted mode:
\begin{equation}
W_{ij}=\max\!\left\{1,\mathrm{round}\!\big(\lambda_w d_{ij}w(i,j)\big)\right\},
\end{equation}
and boundary edges:
\begin{equation}
W_{iB}=\max\!\left\{1,\mathrm{round}\!\big(\lambda_w d_{iB}w(i,i)\big)\right\}.
\end{equation}

\subsection{Why distance behaves like log-likelihood}
If each path segment has effective error probability $p_e<1/2$, then for path length $\ell$:
\begin{equation}
-\log \Prb(\text{path}) \propto \ell\log\frac{1-p_e}{p_e}.
\end{equation}
So minimizing weighted path length is equivalent to maximizing likelihood under this independent-edge approximation.

\subsection{Boundary-augmented matching matrix}
For $k$ defects, LiDMaS+ builds a $2k\times 2k$ matrix:
\begin{align}
C_{ij} &= W_{ij}, && i,j<k,\ i\neq j,\\
C_{i,k+i} &= W_{iB}, && i<k,\\
C_{k+i,k+j} &= 0, && i\neq j.
\end{align}
Solve
\begin{equation}
\min_{M\in\mathcal M}\sum_{(u,v)\in M}C_{uv},
\end{equation}
then lift matched pairs/boundary links to lattice paths and toggle them into $c\in\F^n$.
Consistency check:
\begin{equation}
Hc=s\bmod 2.
\end{equation}

\paragraph{Implementation insight.}
The syndrome-consistency invariant is the primary correctness guardrail: even a low-cost match is invalid if it does not reproduce syndrome parity.

\section{Decoder Objective II: Union-Find + Peeling}
Union-Find decoding is near-linear-time and widely used for topological-code workloads \cite{unionfind2017}. LiDMaS+ follows:
\begin{enumerate}[label=\arabic*.]
  \item start odd clusters at defect vertices,
  \item grow clusters and merge on collisions,
  \item connect unresolved odd roots to boundaries,
  \item build a forest over selected edges,
  \item peel leaf-by-leaf to generate correction edges.
\end{enumerate}

Leaf rule for $e=(v,u)$:
\begin{quote}
if parity$(v)=1$, select $e$ and toggle parity at $u$.
\end{quote}
Output correction is again validated by $Hc=s\bmod 2$.

\section{Reliability Fields and Neural Guidance}
\subsection{Uniform and LLR fields}
Uniform:
\begin{equation}
w(i,j)=1.
\end{equation}
LLR-inspired:
\begin{equation}
w_{\mathrm{LLR}}(p_e)=\max\!\left(0,\log\frac{1-p_e}{p_e}\right),\qquad p_e\in(0,1).
\end{equation}
Implementation maps non-self edges via $p_{\mathrm{data}}$ and boundary/self proxies via $p_{\mathrm{idle}}$.

\paragraph{Interpretation.}
Large $w_{\mathrm{LLR}}$ means ``more confidence in no error,'' so corresponding path choices are penalized more strongly.

\subsection{Neural-guided MWPM plugin}
Feature vector:
\begin{equation}
\phi=(|\Delta x|,|\Delta y|,b),\qquad b\in\{0,1\}.
\end{equation}
Scale model:
\begin{equation}
s(\phi)=\mathrm{clip}\!\left(\beta_0+\beta_1m+\beta_2|\Delta x|+\beta_3|\Delta y|+\beta_4b,\ s_{\min},s_{\max}\right),
\end{equation}
\begin{equation}
\hat W=W\,s(\phi).
\end{equation}
Alternative absolute model:
\begin{equation}
\hat W=\max(0,\alpha m+\beta).
\end{equation}

\paragraph{What this changes mathematically.}
The combinatorial solver stays the same; only edge costs are replaced by model-informed surrogates.

\section{Optional BP Layer (Shared Inference Infrastructure)}
LiDMaS+ also carries BP-style equations used in LDPC/CSS contexts and related simulator modules \cite{stim2021}.

For BSC crossover $p\in(0,1/2)$:
\begin{equation}
L_0=\log\frac{1-p}{p},\qquad
\ell_i^{\mathrm{ch}}=L_0(1-2y_i),
\end{equation}
with erasure override $\ell_i^{\mathrm{ch}}=0$.

Implementation stabilizers:
\begin{equation}
\mathrm{damp}(m_{\mathrm{old}},m_{\mathrm{new}};\eta)=(1-\eta)m_{\mathrm{new}}+\eta m_{\mathrm{old}},
\end{equation}
\begin{equation}
\mathrm{clip}(x;L_{\max})=\max(-L_{\max},\min(x,L_{\max})).
\end{equation}

Check-node update:
\begin{equation}
m_{c\to v}=2\,\mathrm{atanh}\!\left(\prod_{u\in\mathcal N(c)\setminus\{v\}}\tanh\frac{m_{u\to c}}{2}\right),
\end{equation}
with syndrome-sign handling:
\begin{equation}
m_{c\to v}\leftarrow(-1)^{s_c}m_{c\to v}.
\end{equation}
Min-sum approximation:
\begin{equation}
m_{c\to v}=\alpha
\left(\prod_{u\in\mathcal N(c)\setminus\{v\}}\mathrm{sign}(m_{u\to c})\right)
\min_{u\in\mathcal N(c)\setminus\{v\}}|m_{u\to c}|.
\end{equation}
Posterior LLR and decision:
\begin{equation}
\ell_v^{\mathrm{post}}=\ell_v^{\mathrm{ch}}+\sum_{c\in\mathcal N(v)}m_{c\to v},\qquad
\hat e_v=\ind{\ell_v^{\mathrm{post}}<0}.
\end{equation}

\section{Monte Carlo Estimation and Uncertainty}
\subsection{Logical error-rate estimator}
At each sweep point:
\begin{equation}
\widehat{\mathrm{LER}}=\frac{k}{N}.
\end{equation}
This is the sample mean of Bernoulli outcomes $\mathcal L_t$.

\subsection{Wilson 95\% interval}
With $z=1.9599639845$ and $\hat p=k/N$:
\begin{align}
\mathrm{denom}&=1+\frac{z^2}{N},\\
\mathrm{center}&=\frac{\hat p+z^2/(2N)}{\mathrm{denom}},\\
\mathrm{spread}&=\frac{z}{\mathrm{denom}}
\sqrt{\frac{\hat p(1-\hat p)+z^2/(4N)}{N}},
\end{align}
\begin{equation}
\mathrm{CI}_{95\%}=
\left[\max(0,\mathrm{center}-\mathrm{spread}),\min(1,\mathrm{center}+\mathrm{spread})\right].
\end{equation}

\paragraph{Why Wilson instead of normal approximation?}
It behaves better for small $N$ or extreme proportions (near 0 or 1), which are common in high-distance, low-physical-noise regimes.

\subsection{Diagnostics and adaptive stopping}
For any trial statistic $x_t$:
\begin{equation}
\bar x=\frac{1}{N}\sum_{t=1}^N x_t,\qquad
s^2=\frac{1}{N-1}\sum_{t=1}^N(x_t-\bar x)^2,\qquad
\mathrm{SE}(\bar x)=\sqrt{\frac{s^2}{N}}.
\end{equation}
Adaptive stopping (if enabled) terminates when interval half-width $h$ satisfies
\begin{equation}
h\le\varepsilon_{\mathrm{abs}},
\end{equation}
or
\begin{equation}
\frac{h}{\max(\widehat{\mathrm{LER}},10^{-12})}\le\varepsilon_{\mathrm{rel}}.
\end{equation}

\section{Threshold Extraction: Crossings and Scaling}
\subsection{Sweep definitions}
Pauli mode:
\begin{equation}
p\in[p_{\min},p_{\max}],\qquad p\leftarrow p+\Delta p.
\end{equation}
Hybrid mode:
\begin{equation}
\sigma\in[\sigma_{\min},\sigma_{\max}],\qquad \sigma\leftarrow \sigma+\Delta \sigma.
\end{equation}
CSV convention: $\sigma=0$ for Pauli rows, $p=0$ for hybrid rows.

\subsection{Crossing estimators}
Hybrid nearest-grid difference:
\begin{equation}
\sigma_{\mathrm{cross}}\approx\arg\min_{\sigma_i}
\left|
\widehat{\mathrm{LER}}_{d_1}(\sigma_i)-\widehat{\mathrm{LER}}_{d_2}(\sigma_i)
\right|.
\end{equation}
Pauli sign-change interpolation:
\begin{equation}
\Delta(p)=\widehat{\mathrm{LER}}_{d_1}(p)-\widehat{\mathrm{LER}}_{d_2}(p),
\end{equation}
\begin{equation}
p^*=p_k+(p_{k+1}-p_k)\frac{\Delta(p_k)}{\Delta(p_k)-\Delta(p_{k+1})}.
\end{equation}
Bootstrap:
\begin{equation}
k'\sim\mathrm{Binomial}(N,\hat p),\qquad \hat p'=\frac{k'}{N}.
\end{equation}

\subsection{Finite-size scaling collapse}
Using the scaling ansatz often used in threshold studies:
\begin{equation}
x=(p-p_c)d^{1/\nu}.
\end{equation}
LiDMaS+ bins points by $x$ and minimizes
\begin{equation}
\mathcal C(p_c,\nu)=\frac{\sum_b n_b\,\Var(y\mid b)}{\sum_b n_b},
\end{equation}
where $y$ is LER and bins with $n_b<2$ are excluded. Best fit:
\begin{equation}
(\hat p_c,\hat\nu)=\arg\min_{p_c,\nu}\mathcal C(p_c,\nu).
\end{equation}

\section{Deterministic Seeding and Reproducibility}
Each trial seed is a deterministic 64-bit hash of:
\begin{equation}
(\text{base seed},d,\text{sweep key},\text{trial index},\text{thread id}).
\end{equation}
So fixed CLI inputs and thread settings produce reproducible runs, which is central for software-centered research reporting \cite{wayo2026lidmas}.

\section{End-to-End Trial Map}
Each trial executes:
\begin{enumerate}[label=\arabic*.]
  \item sample physical noise (Pauli or CV$\rightarrow$digitized),
  \item compute syndrome $s$,
  \item decode to correction $c$,
  \item verify active-sector invariant $Hc=s\bmod 2$,
  \item evaluate logical indicator $\mathcal L$.
\end{enumerate}
Point estimate:
\begin{equation}
\widehat{\mathrm{LER}}=\frac{1}{N}\sum_{t=1}^N \mathcal L_t.
\end{equation}

\section{Graduate-Level Teaching Flow (Recommended Lecture Plan)}
\begin{enumerate}[label=\arabic*.]
  \item \textbf{Physical context}: surface-code thresholding and why hybrid CV mappings are useful \cite{surfacecode1998,gkp2001,tzitrin2021static,sabo2024weightreduced}.
  \item \textbf{Binary algebra}: derive syndrome equations from parity operators.
  \item \textbf{Decoder objectives}: show MWPM and UF as inference surrogates with different complexity/optimality tradeoffs \cite{pymatching2021,unionfind2017}.
  \item \textbf{Statistical layer}: derive LER estimator, Wilson intervals, and adaptive stopping.
  \item \textbf{Threshold extraction}: crossings, bootstrap uncertainty, and finite-size scaling.
\end{enumerate}

The central message for students is that LiDMaS+ equations form one coherent chain:
\begin{quote}
model assumptions $\rightarrow$ decoder objective $\rightarrow$ logical event definition $\rightarrow$ Monte Carlo estimator $\rightarrow$ threshold inference.
\end{quote}

\bibliographystyle{unsrt}
\bibliography{paper}

\end{document}
